\begin{textsample}{NEG Dim 2 – 2009\_c5 – Score -8.00 – 2009\_c5/enade\_2009\_ci\_cont\_55d90a2674.txt}  \label{ex:f2_neg_017}
Capítulo 5 Distribuição dos Conceitos 5.1 Panorama nacional da distribuição dos conceitos A Tabela 5.1 apresenta a quantidade e o percentual de cursos de CIÊNCIAS CONTÁBEIS participantes do ENADE/2009 que se situaram em cada uma das faixas de conceitos, \textbf{distribuídos} de acordo com a Região em que se localizam os cursos das instituições de educação superior, mantenedoras dos cursos responsáveis pela oferta.
O Gráfico 5.1 apresenta a mesma informação sob a forma de diagrama de caixa para uma melhor visualização da dispersão dos conceitos obtidos.
No diagrama de caixa, a caixa propriamente \textbf{dita} fica limitada pelos quartis.
A barra dentro da caixa em \textbf{negrito} representa a mediana.
Linhas verticais unem os quartis aos valores extremos.
Se estes valores extremos se configurarem como valores \textbf{atípicos}, estes aparecem como pontos isolados e as linhas terminam no maior ( ou menor ) valor ainda não considerado \textbf{atípico}.
Eventuais valores \textbf{atípicos} aparecem como pequenos círculos com a Categoria Administrativa correspondente.
Valores considerados muito \textbf{atípicos} foram representados por estrelas.
Um diagrama de caixa para poucos dados ou com dados muito concentrados pode ter alguns pontos \textbf{coincidentes}.
Por exemplo, se mais de metade dos cursos de uma Região receberem conceito 5, o valor extremo superior, o quartil superior e a mediana coincidirão.
Se o conjunto só contém um dado, este apareceria denotado como a barra da mediana.
Observando -se os dados da Tabela 5.1, percebe -se que, em termos nacionais, \textbf{predominaram} os cursos que receberam conceitos 2 e 3, correspondentes respectivamente a 24,1 \% e a 38,5 \% do total nacional.
Apenas 1,8 \% dos cursos foram classificados com conceito 1 e 3,5 \% dos cursos receberam conceito máximo.
Dos 902 cursos avaliados, 172 ficaram sem conceito, equivalendo a 19,1 \% do total.
Considerando -se separadamente as cinco Regiões, observou -se que, na Região Norte, dos 67 cursos participantes 21 ficaram sem conceito ( 31,3 \% ).
Nenhum curso obteve o conceito máximo e apenas 2 cursos ( 3,0 \% ) obtiveram o conceito mínimo.
Nesta Região, o conceito modal foi 2 ( 37,3 \% ), o conceito 3 foi obtido por 16 cursos ( 23,9 \% ) e o conceito 4 por 3 cursos ( 4,5 \% ) Na Região Nordeste, dos 168 cursos participantes, 26 ficaram sem conceito ( 15,5 \% ).
Obtiveram o conceito máximo apenas 2 cursos ( 1,2 \% ) e 5 cursos ( 3,0 \% ) obtiveram o conceito mínimo.
O conceito modal nesta Região foi o conceito 3, obtido por 65 cursos ( 38,7 \% ) seguido de perto pelo conceito 2, obtido por 54 cursos ( 32,1 \% ).
O conceito 4 foi obtido por 16 cursos ( 9,5 \% ).
Na Região Sudeste, dos 357 cursos participantes, 69 ficaram sem conceito ( 19,3 \% ).
Obtiveram o conceito máximo 15 cursos ( 4,2 \% ) e 3 cursos ( 0,8 \% ) obtiveram o conceito mínimo.
O conceito modal nesta Região foi o conceito 3, obtido por 153 cursos ( 42,9 \% ).
O conceito 2 foi obtido por 70 cursos ( 19,6 \% ) e o conceito 4, por 47 cursos ( 13,2 \% ).
Na Região Sul, dos 197 cursos participantes, 33 ficaram sem conceito ( 16,8 \% ).
Obtiveram o conceito máximo 12 cursos ( 6,1 \% ) e nenhum curso obteve o conceito mínimo.
O conceito modal nesta Região foi o conceito 3, obtido por 83 cursos ( 42,1 \% ).
O conceito 2 foi obtido por 24 cursos ( 12,2 \% ) e o conceito 4, por 45 cursos ( 22,8 \% ).
Na Região Centro-Oeste, dos 113 cursos participantes, 23 ficaram sem conceito ( 20,4 \% ).
Obtiveram o conceito máximo 3 cursos ( 2,7 \% ) e 6 cursos obtiveram o conceito mínimo ( 5,3 \% ).
O conceito modal nesta Região foi o conceito 2, obtido por 44 cursos ( 38,9 \% ).
O conceito 3 foi obtido por 30 cursos ( 26,5 \% ) e o conceito 4, por 7 cursos ( 6,2 \% ). 5.2 Conceitos por Categoria Administrativa e por Região A Tabela 5.2 apresenta a distribuição dos cursos participantes do ENADE/2009 de CIÊNCIAS CONTÁBEIS, por Categoria Administrativa, de acordo com os conceitos por eles alcançados, segundo Região.
O Gráfico 5.2 apresenta a mesma informação sob a forma de diagrama de caixa para uma melhor visualização da dispersão dos conceitos obtidos.
Somente eventuais valores \textbf{atípicos} aparecem com o \textbf{identificador} da Região na qual o curso foi \textbf{habilitado}.
De acordo com as informações apresentadas, 85,7 \% dos cursos eram vinculados a instituições privadas, 5,4 \% a instituições federais, 5,4 \% a instituições estaduais e 3,4 \% a instituições municipais.
As instituições federais foram as que apresentaram, em termos nacionais, maior proporção de cursos com conceito 5, equivalentes a 26,5 \% do total dessa categoria.
O conceito modal nas instituições federais foi o conceito 4 ( 36,7 \% ) seguido do conceito 3 ( 26,5 \% ), sendo que nenhum curso obteve o conceito mínimo ou o conceito 2.
Dos 49 cursos avaliados em instituições federais, 5 ( 10,2 \% ) ficaram sem conceito.
Entre as instituições estaduais, o conceito modal foi o conceito 3 ( 36,7 \% ), sendo que nenhum curso obteve o conceito mínimo e 12,2 \% o conceito máximo.
Dos 49 cursos avaliados em instituições estaduais, 1 curso ( 2,0 \% ) ficou sem conceito.
O conceito modal nas instituições municipais foi o conceito 3 ( 51,6 \% ) seguido do conceito 2 ( 16,1 \% ), sendo que nenhum curso obteve o conceito máximo e 5 cursos ( 16,1 \% ) obtiveram o conceito 2.
O conceito mínimo foi atribuído a 2 cursos ( 6,5 \% ).
Dos 31 cursos avaliados em instituições municipais, 3 ( 9,7 \% ) ficaram sem conceito.
Entre as instituições privadas, o conceito modal foi o conceito 3 ( 38,8 \% ), seguido do conceito 2 ( 26,1 \% ), sendo que 1,8 \% dos cursos obtiveram o conceito mínimo e 1,7 \% o conceito máximo.
Dos 773 cursos avaliados em instituições privadas, 163 ( 21,1 \% ) ficaram sem conceito.
Considerando -se a análise por Região observou -se que, em todas as Regiões, cursos de todas as Categorias Administrativas consideradas se fazem presentes sendo que as instituições privadas representam uma grande maioria dos cursos.
Dos 67 cursos avaliados na Região Norte, 21 ficaram sem conceito sendo 15 em instituição federal ( 12,5 \% ) e 20 em instituições privadas ( 35,1 \% ).
O conceito modal nas instituições federais foi o conceito 3 ( 50,0 \% ), seguido do conceito 4 ( 37,4 \% ).
O único curso com conceito em instituição estadual obteve o conceito 2 e o único curso em instituição municipal participante obteve o conceito mínimo.
O conceito modal nas instituições privadas foi 2 ( 42,1 \% ), seguido do conceito 3 ( 21,1 \% ).
Nenhuma instituição na Região obteve o conceito máximo, sendo que o conceito mínimo foi atribuído apenas a instituições privadas e à única municipal já mencionada.
Na Região Nordeste, dos 168 cursos participantes, 26 ficaram sem conceito sendo 2 em instituições federais ( 13,3 \% ), 1 em instituição estadual ( 5,3 \% ), 2 em instituições municipais ( 66,7 \% ) e 21 em instituições privadas ( 16,0 \% ).
O conceito modal nas instituições federais foi o conceito 4 ( 53,3 \% ), seguido do conceito 3 ( 33,3 \% ).
Nas instituições estaduais o conceito modal foi 3 ( 47,4 \% ), seguido do conceito 2 ( 26,3 \% ) e do conceito 4 ( 15,8 \% ).
O único curso avaliado em instituição municipal obteve o conceito 3.
O conceito modal nas instituições privadas foi 3 ( 38,2 \% ), seguido do conceito 2 ( 37,4 \% ).
Apenas 1 curso em instituição estadual e 1 curso em instituição privada na Região obtiveram o conceito máximo.
O conceito mínimo foi atribuído apenas a 5 cursos em instituições privadas.
Na Região Sudeste, dos 357 cursos participantes, 69 ficaram sem conceito sendo 2 em instituições federais ( 15,4 \% ) e 67 em instituições privadas ( 20,2 \% ).
Os conceitos modais nas instituições federais foram os conceitos 4 e 5 ( 38,5 \% cada ).
Nas instituições estaduais os conceitos 3, 4 e 5 foram obtidos por 1 curso cada.
Nas instituições municipais o conceito modal foi o conceito 3 ( 80,0 \% ), seguido do conceito 2 ( 20,0 \% ).
O conceito modal nas instituições privadas foi 3 ( 43,2 \% ), seguido do conceito 2 ( 20,5 \% ) e do conceito 4 ( 12,4 \% ).
Apenas 1 curso em instituição estadual, 5 cursos em instituições federais e 9 cursos em instituições privadas na Região obtiveram o conceito máximo.
O conceito mínimo foi atribuído a apenas 3 cursos de instituições privadas ( 0,9 \% ).
Na Região Sul, dos 197 cursos participantes, 33 ficaram sem conceito todos em instituições privadas.
Os 6 cursos de instituições federais participantes obtiveram o conceito 5.
Nas instituições estaduais o conceito modal foi o conceito 4 ( 38,9 \% ), seguido dos conceitos 3 ( 27,8 \% ) e 5 ( 22,2 \% ).
Nas instituições municipais o conceito modal foi o conceito 3 ( 50,0 \% ), seguido do conceito 4 ( 41,7 \% ).
O conceito modal nas instituições privadas foi 3 ( 44,7 \% ), seguido do conceito 3 ( 20,5 \% ) e do conceito 2 ( 13,0 \% ).
O conceito máximo foi atribuído a 6 cursos em instituições federais, 4 em instituições estaduais e 2 em instituições privadas.
Nenhum curso na Região obteve o conceito mínimo.
Na Região Centro-Oeste, dos 113 cursos participantes, 23 ficaram sem conceito sendo 1 em instituição municipal e 22 em instituições privadas.
Os 7 cursos de instituições federais participantes obtiveram os conceitos 3 ( 42,9 \% ), 4 ( 28,6 \% ) e 5 ( 28,6 \% ).
Nas instituições estaduais os conceitos modais foram os conceitos 3 e 4 ( 37,5 \% cada ), seguidos do conceito 2 ( 25,0 \% ).
Nas instituições municipais o conceito modal foi o conceito 2 ( 40,0 \% ), seguido dos conceitos 1 e 3 ( 20,0 \% cada ).
O conceito modal nas instituições privadas foi 2 ( 43,0 \% ), seguido do conceito 3 ( 24,7 \% ).
O conceito máximo foi atribuído a 3 cursos, sendo 2 em instituições federais e 1 em instituições privadas.
O conceito mínimo foi atribuído a 1 curso em instituição municipal e a 5 cursos em instituições privadas. 5.3 Conceitos por Organização Acadêmica e por Região Na Tabela 5.3 encontra -se a distribuição dos conceitos atribuídos aos cursos participantes do ENADE/2009 de CIÊNCIAS CONTÁBEIS, por Organização Acadêmica, segundo as Regiões.
O Gráfico 5.3 apresenta a mesma informação sob a forma de diagrama de caixa para uma melhor visualização da dispersão dos conceitos obtidos.
Somente eventuais valores \textbf{atípicos} aparecem com o \textbf{identificador} da Categoria Administrativa.
Os cursos desta área estão representados em Universidades, Centros Universitários e Faculdades.
Não houve participação de nenhum curso em Centro Federal de Educação Tecnológica ou em IFECT.
Percebe -se que, em termos nacionais, \textbf{predominaram} os cursos que receberam conceitos 3, 2 e 4, correspondentes respectivamente a 347, 217 e 118 cursos do total nacional de 902 cursos.
Deste total de 902 cursos, 172 ficaram sem conceito, sendo 34 em Universidades, 7 em Centros Universitários e 131 em Faculdades.
O conceito modal nas Universidades, Centros Universitários e Faculdades foi o conceito 3, seguido do conceito 2, nos Centros Universitários e Faculdades, e do conceito 4 nas Universidades.
Considerando -se separadamente as cinco Regiões, observa -se que, na Região Norte, dos 67 cursos participantes 25 obtiveram conceito 2, 16 o conceito 3 e 21 ficaram sem conceito, sendo 2 em Universidades e 19 em Faculdades.
Nenhum curso obteve o conceito máximo e apenas 2 cursos, 1 em Centro Universitário e 1 em Faculdade, obtiveram o conceito mínimo.
O conceito modal nas Universidades e Centros Universitários foi o conceito 3 e nas Faculdades foi o conceito 2.
Na Região Nordeste, dos 168 cursos participantes 65 obtiveram conceito 3, 54 o conceito 2 e 26 ficaram sem conceito, sendo 5 em Universidades e 21 em Faculdades.
Apenas 2 cursos, 1 em Universidade e outro em Faculdade obtiveram o conceito máximo.
Obtiveram o conceito mínimo 5 cursos, todos em Faculdades.
O conceito modal nas Universidades foi o conceito 3, seguido do conceito 4, nos Centros Universitários foi o conceito 2 seguido do 3 e nas Faculdades foi o conceito 3, seguido do conceito 2.
Na Região Sudeste, dos 357 cursos participantes 153 obtiveram conceito 3, 70 o conceito 2 e 47 o conceito 4.
Ficaram sem conceito 69 cursos, sendo 16 em Universidades, 7 em Centros Universitários e 46 em Faculdades.
Obtiveram o conceito máximo 15 cursos, 10 em Universidades, 2 em Centros Universitários e 3 em Faculdades.
Obtiveram o conceito mínimo apenas 3 cursos em Faculdades.
O conceito modal nas Universidades, Centros Universitários e Faculdades foi o conceito 3, seguido do conceito 4 nas Universidades e do conceito 2, nos Centros Universitários e Faculdades.
Na Região Sul, dos 197 cursos participantes 83 obtiveram conceito 3, 45 o conceito 4 e 24 o conceito 2.
Ficaram sem conceito 33 cursos, sendo 8 em Universidades e 25 em Faculdades.
Obtiveram o conceito máximo 12 cursos, todos em Universidades.
Nenhum curso obteve o conceito mínimo.
O conceito modal nas Universidades, Centros Universitários e Faculdades foi o conceito 3, seguido do conceito 4 nas Universidades e Centros Universitários, e dos conceitos 2 e 4, nas Faculdades.
Na Região Centro-Oeste, dos 113 cursos participantes 44 obtiveram conceito 2, 30 o conceito 3 e 7 o conceito 4.
Ficaram sem conceito 23 cursos, sendo 3 em Universidades e 20 em Faculdades.
Obtiveram o conceito máximo 3 cursos, sendo 2 em Universidades e 1 em Faculdade.
O conceito mínimo foi obtido por 6 cursos, sendo 2 em Universidades e 4 em Faculdades.
O conceito modal nas Universidades foi o conceito 3, seguido do conceito 2.
Nos Centros Universitários e Faculdades foi o conceito 2, seguido do conceito 3.

% matched lemmas: atípico, coincidente, distribuir, dito, habilitado, identificador, negrito, predominar
\end{textsample}
